\documentclass{article}
\usepackage[T1]{fontenc}
\usepackage{lmodern}
\usepackage{graphicx}
\usepackage{amsmath,amssymb}
\usepackage[a4paper,margin=2cm]{geometry}
\usepackage[absolute,overlay]{textpos}
\setlength{\TPHorizModule}{1mm}
\setlength{\TPVertModule}{1mm}
\usepackage[indonesian]{babel}
\usepackage{hyperref}
\setlength{\emergencystretch}{2em}
% unit: Halaman 17

\begin{document}

% === Halaman 17 (python/native) ===

Serangan pada Diffie-Hellman

• Pertukaran kunci Diffie-Hellman hanya aman terhadap serangan pasif, misalnya 

penyadapan (eavesdropping).

• Eva yang menyadap komunikasi antara Alice dan Bob hanya mengetahui nilai g, p, A dan B 

(semuanya publik). Dia tidak dapat mendeduksi a dan b (kunci privat Alice dan Bob) dari 

nilai-nilai publik tersebut. Untuk mendeduksi a dan b, Eva harus dapat menghitung 

logaritma diskrit dari persamaan A = ga mod p atau B = gb mod p.

• Namun, Diffie-Hellman tidak tahan terhadap serangan aktif seperti man-in-the-middle 

attack. 

• Misalkan Mallory mengintervensi komunikasi antara Alice dan Bob. Kepada Alice Malllory 

menyamar (impersonation) sebagai Bob, dan kepada Bob Mallory menyamar sebagai Alice.

Rinaldi M/II4020 Kriptografi

17

\end{document}
